\chapter{Game-Theoretic Framework}\label{ch:game_theory}

The sign-separated streams of Chapter~\ref{ch:decompositions} expose accumulated opposing contributions. This chapter formalizes the backward-game picture of Figure~\ref{fig:game_trajectory_overview}, tracing backward rules through these contributions. The stopping decomposition supplies a common state graph for two games whose policies represent gradients and relevance propagation. Their trajectory distributions support method design (Chapter~\ref{ch:saliency}) and diagnostics (Chapter~\ref{ch:trajectory}). Exact recovery concerns the anchor rules; myopic temperature changes and risk-sensitive policies deform these rules. The companion supplement, Chapter~\ref{app:supp_games}, gives detailed definitions and proofs.

\section{Related Work}\label{ptnb:sec:related-work}\label{sec:related_work_games}

\textcite{Ancona2018} unify several gradient-based rules as modified gradients, a map-level construction excluding the $\alpha\beta$ rule. Our signed trajectory-space representation unifies gradients and the $\alpha\beta$ rule through their routing laws. Together with the split representation, it distinguishes which internal signals attribution retains and how it propagates them backward.

\textbf{Game and trajectory views.} The closest mathematical precursor, the ReLU Net Game of
\textcite{Gaubert2025}, identifies forward ReLU-network evaluation with a backward Stopping Game's
value, without connecting it to the backward pass or attribution. We reformulate it through a novel Stopping
Decomposition whose weight split relates to the \ac{dc} decomposition used directly for attribution
by \textcite{hiddenMono}. Our \ac{sg} and \ac{rg} use this weight split to recover backward attribution
rules rather than the pure DC decomposition.
The \ac{nrm} framework of \textcite{Xiong2026nrm} assigns normalized signed relevance to network walks.
Section~\ref{sec:nrm_comparison} compares its treatment of \ac{lrp} and concept relevance propagation with our
trajectory probabilities and discounted occupations.

\textbf{Relevance propagation and stochastic explanations.} \ac{lrp}, Deep Taylor, and Excitation
Backprop form the core relevance propagation family~\parencite{Bach2015,Montavon2017DeepTaylor,Zhang2018pointinggame}; variants modify redistribution for contrast, stability, or
signed relevance~\parencite{Binder2016LRN,Kohlbrenner2020,Gu2019CLRP,Nam2020,vukadin2024abslrp}. Our common
trajectory-space semantics recovers gradients through \ac{sg} and the generalized $\alpha\beta$ rule
through \ac{rg}, whose analytical modes give principled modifications of backward rules.
Excitation Backprop~\parencite{Zhang2018pointinggame} is an absorbing Markov chain whose redistribution
on non-negative activations gives the unstabilized $z^+$-rule with biases omitted: the
$(\alpha,\beta)=(1,0)$ case of $\alpha\beta$-\ac{lrp}. Our game adds an opponent for negative-weight
influence and replaces input-dependent transitions with soft Bellman equations.
Deep Taylor~\parencite{Montavon2017DeepTaylor} justifies $z^+$ on non-negative inputs;
\ac{rg} constructs a game for the generalized $\alpha\beta$ rule.

\textbf{Shapley-value attribution methods.} KernelSHAP explains marginal input-feature contributions
relative to a baseline through a cooperative \emph{black-box} game over feature
coalitions~\parencite{LundbergLee2017}. Our complementary \emph{white-box} game traces relevance
through network internals.

\textbf{Robust explanations via Softplus.} \textcite{Dombrowski2020} propose Softplus replacement as a
curvature-based defense against gradient-explanation manipulation.
Section~\ref{ptnb:app:softplus-stopping-game} shows that entropy-regularizing the activation-state
stop/continue choice in \ac{sg} recovers Softplus-network gradients, situating their construction
within our framework.

\textbf{Interpretability in Transformers and ViTs.} Attention Rollout / Flow use attention as a routing
graph~\parencite{Abnar2020attention}; transformer-\ac{lrp} propagates through attention and skip
connections~\parencite{Chefer2021GenericAttention,Ali2022AttnLRP,Achtibat2024AttnLRP}; DAVE addresses ViT
tokenization artifacts through conditioned operators and equivariant averaging~\parencite{Wrobel2024}.
Our game semantics represents conditioned operators through a QK-softmax oracle reference policy
and Value-routing attribution trajectories. Its \ac{rg}+Eq mode uses Reynolds averaging to suppress
DAVE's targeted tokenization-grid artifacts, with further control through the game modes.
\section{Backward Attribution}\label{ptnb:sec:background}

Throughout the game, attention, and trajectory derivations, $W^{(l)}_{ij}$ denotes the edge from
predecessor $i$ to successor $j$: the adjacency matrices transpose the DC recurrences' output-by-input
matrices while representing the same network. The tag $\mathrm{stop}$ distinguishes stopping from
DC streams. Consider a ReLU network with layers $l = 1,\dots,L$, input $a^{(0)} = x$, pre-activations
$z^{(l)}_j = \sum_i a^{(l-1)}_i W^{(l)}_{ij} + b^{(l)}_j$, and activations $a^{(l)}_j =
    \relu(z^{(l)}_j)$. We initially explain a scalar activation $f(x) = a^{(L)}_u$;
Section~\ref{ptnb:app:affine-readout} extends this to ungated, possibly negative affine logits.
Write $[u]^{+} = \max(u,0)$, $[u]^{-} = \max(-u,0)$, hence $u = [u]^+ - [u]^-$, and
$W^{(l,\pm)} = [W^{(l)}]^{\pm}$ entrywise.

\paragraph{Layer-wise relevance propagation.} Starting from $R^{(L)}_u = f(x)$,
\ac{lrp}~\parencite{Bach2015} redistributes each neuron's relevance $R^{(l)}_j$ backward, layer by layer, in proportion to each input's
pre-activation contribution $a^{(l-1)}_i W^{(l)}_{ij}$, with bias and stabilizer shares absorbed.
Rules weight these contributions differently. The $\alpha\beta$-rule uses $\alpha,\beta\geq0$;
$\alpha-\beta=1$ gives conservation when both participating sign streams are nonempty and bias,
stabilizer, and empty-stream absorption are accounted for. With $\varepsilon\geq0$
\ch{added to each sign denominator, as in common implementations~\parencite{Anders2021Zennit}},
positive and negative fractions are normalized separately and scaled by $\alpha$ and $\beta$:
\begin{align}\label{ptnb:eq:lrp-ab}
    R^{(l-1)}_i &= \textstyle\sum_j \Bigl( \alpha\, \tfrac{[a^{(l-1)}_i
    W^{(l)}_{ij}]^+}{\Sigma^{(l,+)}_j + \varepsilon} - \beta\, \tfrac{[a^{(l-1)}_i
    W^{(l)}_{ij}]^-}{\Sigma^{(l,-)}_j + \varepsilon} \Bigr) R^{(l)}_j, \nonumber\\
    \Sigma^{(l,\pm)}_j &= \textstyle\sum_k [a^{(l-1)}_k W^{(l)}_{kj}]^{\pm} + [b^{(l)}_j]^{\pm}.
\end{align}
Relevance assigned to the bias and to $\varepsilon$ is absorbed. For non-negative activations,
$(\alpha,\beta)=(1,0)$ with $\varepsilon=0$ and biases omitted from normalization gives the
$z^+$-rule of Deep Taylor decomposition~\parencite{Montavon2017DeepTaylor} and
Excitation Backprop~\parencite{Zhang2018pointinggame}; other rules (LRP-$\varepsilon$, LRP-$\gamma$)
change the redistribution fractions.

\paragraph{Transformers.} In a Vision Transformer (\acs{vit}), \ch{LayerNorm couples the features of a token through its
input-dependent mean and variance,} and the attention matrix $A = \softmax(QK^\top/\sqrt{d})$
depends on the input, which breaks conservation~\parencite{Ali2022AttnLRP}.
\ch{Transformer-specific rules therefore \emph{condition} on the forward pass: conservative
propagation~\parencite{Ali2022AttnLRP} and DAVE~\parencite{Wrobel2024} hold $A$ and the LayerNorm denominator fixed,
retain linear mean-centering, and propagate attention relevance through the Value path only;
AttnLRP~\parencite{Achtibat2024AttnLRP} also propagates through the softmax and the QK product.}

These rules route mass backward from each neuron to its inputs; the games below make that routing explicit.
\section{Attribution as a Game}\label{ptnb:sec:game}

\paragraph{Setup and common state-space structure.}
Using Section~\ref{ptnb:sec:background}'s notation, we construct a finite-horizon two-player
game~\parencite{OsborneRubinstein1994} played \emph{backward} from the output. States record layer,
neuron, and game-specific player-turn and stream-sign labels. At each non-terminal state, the active
player stops at an absorbing cemetery state~$\dagger$ or continues. Skip connections randomly select
summand states; max-pooling greedily selects maximal-payoff states.
Section~\ref{ptnb:app:notation} gives the labels (Table~\ref{ptnb:tab:notation}).
For player $p$ and opponent $p'$, the \ac{sg} value $V_p(s)$ is the $\gamma$-discounted expected
payoff difference $r_p-r_{p'}$ along trajectories from $s$ under the policies. For $p$ in turn at $s$,
the Bellman recursion is
\begin{align}\label{ptnb:eq:bellman}
    V_p(s) = \max_{a \in \mathcal{A}_x(s)} \Big\{ &r_p(s,a) - r_{p'}(s,a) \nonumber\\
    &+ \gamma(s) \textstyle\sum_{s'} \P_x(s' \mid s,a)\, V_p(s') \Big\},
\end{align}
and the \acl{rg} instead entropy-regularizes mixed routing decisions, yielding the soft game
value~\eqref{ptnb:eq:soft-bellman}.
Both games start with unit mass at the positive output state.
Player policies induce a backward trajectory measure $\Pi_x$ on $\zeta=(\zeta_0,
    \zeta_1, \dots)$. With game-specific cumulative discount $d_t(\zeta)$ up to time $t$, define the
\textbf{discounted occupation measure}
\begin{align}\label{ptnb:eq:attribution}
    \Gamma_x(v) \coloneqq \E_{\zeta \sim \Pi_x}\!\left[ \textstyle\sum_t d_t(\zeta)\,\1\{\zeta_t = v\}
    \right].
\end{align}
It measures the discounted weight visiting each state.

\subsection{The Stopping Game}\label{ptnb:sec:stopping}\label{sec:stopping_game_reinterpretation}

Each layer-$l$ neuron $j$ has states $s^{(l)}_{j,+}$ and $s^{(l)}_{j,-}$; the sign identifies the player in turn.
Continuation pays the bias and follows weight-magnitude environment transitions with discount $\gamma_j^{(l)} = \sum_i |W_{ij}^{(l)}|$ (positive- and
negative-weight predecessors follow $\P \propto W^{(l,\pm)}/\gamma^{(l)}_j$, and a negative weight
passes the turn to the opponent).
If $\gamma_j^{(l)}=0$, continuation pays the bias and enters the cemetery. We select $\operatorname{ReLU}'(0)=0$ and the smallest-index maximizer at pooling ties.
The expected payoffs realize the novel Stopping Decomposition into positive and negative paths to
each activation (Lemma~\ref{ptnb:lem:parity-path}); the gradient is the difference of two
\emph{non-negative} occupation measures:

\begin{theoremenv}[Gradient representation, \LeanResult{ptnb:thm:gradient}]\label{ptnb:thm:gradient}
The Stopping Game allows a subgame-perfect Nash equilibrium, unique up to tie-breaking at
zero-activation neurons.
On strict gate cells (and with unique max-pooling winners where applicable), under the stop-at-tie-breaks equilibrium the discounted occupation measure $\Gamma_x(s_{i,p}^{(l)})
    \ge 0$ equals the output derivative with respect to the corresponding \emph{Stopping Decomposition}
activation. At ties, the corresponding identities refer to the backward quantities selected by the stated gate and winner conventions.
On these cells, the gradient of $f$ with respect to the input is the difference $\Gamma_x(s_{i,+}^{(0)}) -
    \Gamma_x(s_{i,-}^{(0)})$.\end{theoremenv}
Section~\ref{ptnb:app:relu-net-game} supplies the proof, an example, the full skip/pooling state
space, and a comparison with the \ac{dc} decomposition.

Modeling $\alpha\beta$-\ac{lrp} or attention modules via the \ac{sg} would require
activation-dependent transitions. The \ac{rg} moves activation-dependence into the players'
policies, opening a continuous family of deformations.

\subsection{The Routing Game}\label{ptnb:sec:routing}\label{subsec:routing_game_def}

Each \ac{rg} neuron has an activation state $s^{(l,\mathrm{act})}_{j,q}$ and linear routing states
$s^{(l,\mathrm{lin})}_{j,q,+}$ and $s^{(l,\mathrm{lin})}_{j,q,-}$ for the weight parts $W^{(l,+)}$
and $W^{(l,-)}$. At the activation state, the player stops or continues; on continuation, each stream
is chosen with probability $\tfrac12$; the $-$ stream switches players, and \ch{game-value discount $2$}
compensates the split.
At linear state $s^{(l,\mathrm{lin})}_{j,q,\kappa}$ for layer-$l$ neuron $j$ and sign stream
$\kappa\in\{+,-\}$, the player $q$ in turn mixes over predecessors $i$, a bias
node, and an outside option $\perp_{\mathrm{lin}}$ under policy
$\pi_{q}(s_{j,q,\kappa}^{(l,\text{lin})})$. The latter two absorb the trajectory with log-payoffs
$\log[b^{(l)}_j]^\kappa$ and $\log\varepsilon$.
In log space, player $r$'s advantage of moving to predecessor $i$ becomes the additive action--value
\begin{align}\label{ptnb:eq:q-value}
    Q_r\bigl(s^{(l,\mathrm{lin})}_{j,q,\kappa}, i\bigr) = \begin{cases} \log W_{ij}^{(l,\kappa)} {+} V_r(s^{(l-1,\mathrm{act})}_{i,q}), & r{=}q,\, l{\ge}2,\\[0.2em] \log \bigl[x_i W_{ij}^{(1)}\bigr]^{\kappa}, & r{=}q,\, l{=}1, \\
        -\infty, & r \neq q,
    \end{cases}
\end{align}
where $V_r$ is the successor activation state's \emph{game value} for player $r$. Shannon-entropy
regularization of each linear routing subgame~\parencite{Geist2019regularized} yields the equilibrium
value and Gibbs best response:
\begin{align}
    V_q\!\bigl(s^{(l,\mathrm{lin})}_{j,q,\kappa}\bigr) &= \tau \log \textstyle\sum_i
    \exp\!\bigl(Q_q(s^{(l,\mathrm{lin})}_{j,q,\kappa}, i)/\tau\bigr), \label{ptnb:eq:soft-bellman}\\
    \pi_{q}(s_{j,q,\kappa}^{(l,\text{lin})})(i) &\propto
    \exp\!\bigl(Q_q(s^{(l,\mathrm{lin})}_{j,q,\kappa}, i)/\tau\bigr), \label{ptnb:eq:gibbs-policy}
\end{align}
where the sums include the bias node and $\perp_{\mathrm{lin}}$. For $\tau>0$, an SPNE exists;
at $\tau=1$, its occupation measure represents $\alpha\beta$-\ac{lrp}:

The Routing Game cumulative trajectory discount multiplies by $2\alpha$ on entering a positive linear-routing state, $2\beta$ on entering a negative one, and $1$ otherwise, separately from game-value discounting.

\begin{theoremenv}[$\alpha\beta$-\ac{lrp} representation, \LeanResult{ptnb:thm:lrp-recovery}]\label{ptnb:thm:lrp-recovery}\label{thm:lrp_recovery}
For $\tau>0$, the Gibbs policies of Eq.~\eqref{ptnb:eq:gibbs-policy} and stopping decisions from the corresponding continuation values yield a canonical SPNE of the entropy-regularized \ac{rg}. We select stopping at every activation log tie and absorption for an empty routing stream, as specified in Remark~\ref{ptnb:rem:routing-spne-uniqueness}. At the anchor $\tau=1$, the selected gate is $\mathbf{1}[z>0]$ and the
$\alpha\beta$-\ac{lrp} attribution map in~\eqref{ptnb:eq:lrp-ab} equals $R_i^{(l)} =
    f(x)\,\bigl(\Gamma_x(s_{i,+}^{(l)}) - \Gamma_x(s_{i,-}^{(l)})\bigr)$.\end{theoremenv}
Section~\ref{ptnb:app:routing-formal} gives the full game definition, the proofs, and a worked example
of the local routing dynamics (Figure~\ref{ptnb:fig:Routing-Game-illustration}).

In practice, $\tau$ is a \emph{myopic} control: each linear subgame uses its $\tau$-scaled Gibbs
policy on its own action values, with downstream game values fixed at $\tau=1$, preserving the network's actual forward
activations. Re-solving downstream at $\tau$ would change them. Thus $\tau\neq1$ deforms the
occupation measure; low $\tau$ concentrates routing on dominant pathways.

\section{Game-Theoretic Modes}\label{ptnb:sec:modes}\label{sec:risk_aversion_regularization}\label{ptnb:sec:controls}

Risk aversion and entropy regularization yield novel backward-rule modes for probing and tuning
explanations (Table~\ref{ptnb:tab:modes}). Section~\ref{ptnb:sec:game}'s exact representations are the
anchors; modes beyond the \ac{lrp} parameters $(\alpha,\beta,\varepsilon)$ are deformations, inactive at the anchor.
\ch{We use five: temperature $\tau$ (focus), activation-state risk aversion $\lambda$ and its
attention-logit variant $\lambda_{\mathrm{sm}}$, the probit gate $\tilde\Phi$ (smoothing), and equivariant averaging
over $N$ input transformations (grid artifacts).}

\ch{\emph{Focus.} The myopic temperature of Section~\ref{ptnb:sec:routing} rescales the Gibbs policy at every linear state while keeping the $\tau{=}1$ game values. Low $\tau$ concentrates routing on dominant predecessors; high $\tau$ spreads mass more uniformly, exposing broader contextual features as well as architectural priors.}
\ch{\emph{Risk aversion.} \Acl{ptnbADF} (\acs{ptnbADF};~\parencite{gast2018lightweight}) propagates a
Gaussian mean $\mu$ and variance $\sigma^2$ for every pre-activation $z$ from an input variance
$\sigma^2$ in one additional moment pass. At activation states, the risk-averse player replaces $z$ by the lower-confidence bound
$\mu+\lambda\sigma$, $\lambda<0$, attenuating high-variance routes while leaving linear-state routing
unchanged. This is the pessimistic counterpart of confidence-bound
exploration~\parencite{auer2002finite,zhang2021risk}. On attention
blocks, $\lambda_{\mathrm{sm}}$ applies the same shift to the oracle's attention logits before the
softmax, downweighting high-variance key tokens.}
\ch{\emph{Gate smoothing.} The risk-averse stop/continue rule is the hard gate
$\1[\mu/\sigma+\lambda>0]$. The probit gate $\tilde\Phi$ lets the player instead continue with
probability $\Phi(\mu/\sigma+\lambda)$, the marginal of a player who perceives the continuation
payoff as $z+\eta$, $\eta\sim\mathcal N(0,1)$, and continues iff it is positive. At $\mu{=}z$,
$\sigma{=}1$, $\lambda{=}0$ this probit-regularized choice recovers the Gaussian-smoothed ReLU
$\E[\relu(z+\eta)]$, just as entropy regularization recovers
Softplus (Section~\ref{ptnb:app:softplus-stopping-game}); on ReLU networks it is a smoothing
deformation. Because the gate uses the \ac{ptnbADF} $\sigma$, $\sigma^2$ is active whenever
$\tilde\Phi$ is on.}
\ch{\emph{Exact modes and cost.} $(\alpha,\beta)$ weight the two streams of the occupation
measure, and the \ac{lrp} stabilizer $\varepsilon$ adds an outside option $\perp_{\mathrm{lin}}$
with log-score $\log\varepsilon$ at every linear state, leaking weak-evidence neurons into the
cemetery while leaving strong-evidence ones near-unaffected (Theorem~\ref{ptnb:thm:ab-recovery}). The
exact modes and $\tau$ add no cost; $\lambda$, $\lambda_{\mathrm{sm}}$ and $\tilde\Phi$ add one
\ac{ptnbADF} pass ($\vitRGanchorTime\to\vitRGlocTime$\,s per ViT-B/16 image); averaging multiplies the
cost by $N$ ($4.93$\,s at $N{=}50$, $\vitRGEqTime$\,s with $\lambda_{\mathrm{sm}}$).
Section~\ref{ptnb:app:modes} derives the modes and gives defaults, costs and an order of use.}

\begin{table}[t]
    \centering
    \caption[Modes of the games]{Modes of the games. ``Exact'': the mode stays within a represented rule; ``deform.'':
    optional deformation, off (default value) at the anchor. $\sigma^2$: input variance of the \ac{ptnbADF}
    moment propagation; $N$: number of transformed inputs averaged by \ac{rg}+Eq.}\label{ptnb:tab:modes}
\TableContent{
    \small
    \setlength{\tabcolsep}{4pt}
    \begin{tabularx}{\textwidth}{@{}>{\raggedright\arraybackslash}p{0.08\textwidth}>{\raggedright\arraybackslash}X>{\raggedright\arraybackslash}p{0.08\textwidth}>{\raggedright\arraybackslash}p{0.24\textwidth}@{}}
        \toprule
        \textbf{Mode} & \textbf{Effect on the game} & \textbf{Default} & \textbf{Property} \\
        \midrule
        $(\alpha,\beta)$ & branch discounts of the two streams (exact) & $(2,1)$ & support / oppose \\
        $\varepsilon$ & outside option $\perp_{\mathrm{lin}}$, payoff $\log\varepsilon$ (exact) & $0.5$ &
        damps low-inflow neurons \\
        $\tau$ & entropy regularization of routing (deform.) & $1$ & sharp / diffuse \\
        $\lambda$, $\sigma^2$ & values $\to$ lower-confidence bounds (deform.) & $\lambda{=}0$ & noise /
        fragility \\
        $\lambda_{\mathrm{sm}}$ & attention logits shifted by $\lambda_{\mathrm{sm}}$ times a per-key
        \ac{ptnbADF} std (deform.) & $0$ & stable attention \\
        $\tilde\Phi$ & noise on player observations (deform.) & off & gate smoothing \\
        $N$ & average over input transformations (deform.) & off & grid artifacts \\
        \bottomrule
    \end{tabularx}
}%
\end{table}


\section{Signed affine output logits}\label{ptnb:app:affine-readout}

The activation-output formulation above extends to the signed logits of a classifier by an
\emph{ungated affine readout}. Let $h_j=a_j^{(L)}\geq0$ be the last hidden activations and
\begin{align}\label{ptnb:eq:affine-readout}
    f_u(x)=\sum_j c_{uj}h_j(x)+b_u.
\end{align}
This output has no ReLU or sign-dependent stopping gate.

\paragraph{Stopping Game.}
Add a root with forced continuation, immediate bias payoff difference $b_u$, and structural
discount $\gamma_u=\sum_j|c_{uj}|$. When $\gamma_u>0$, it moves to hidden neuron $j$ with
probability $|c_{uj}|/\gamma_u$, retaining the player for $c_{uj}\geq0$ and switching players
for $c_{uj}<0$. Its value is therefore $b_u+\sum_jc_{uj}h_j=f_u(x)$, even when this value is
negative. If $\gamma_u=0$, the root pays $b_u$ and moves to the cemetery; the input derivative
is zero. On the strict gate cells of Theorem~\ref{ptnb:thm:gradient}, let
$\Gamma^{\mathrm{SG}}_{j,\pm}(i)$ denote the input-state occupations from a unit source at
hidden neuron $j$ with initial player $+$. The added transition gives
\begin{align}\label{ptnb:eq:affine-readout-gradient}
    \partial_i f_u(x)=\sum_j c_{uj}
    \bigl(\Gamma^{\mathrm{SG}}_{j,+}(i)-\Gamma^{\mathrm{SG}}_{j,-}(i)\bigr).
\end{align}
This chain rule retains non-negative path occupations and the selected-backward convention at gate
or pooling ties.

\paragraph{Routing Game.}
Append an unconditionally continuing, ungated source to the hidden Routing Game, with positive and negative
branch masses including the bias:
$\Sigma_u^\kappa=\sum_j[h_jc_{uj}]^\kappa+[b_u]^\kappa$.
At $\tau=1$, the sign-oracle probabilities $1/2$, trajectory discounts $2\alpha$ and $2\beta$,
and the branchwise Gibbs policies give the coefficients
\begin{align}\label{ptnb:eq:affine-readout-routing}
    q_j^+=\alpha\frac{[h_jc_{uj}]^+}{\Sigma_u^++\varepsilon},
    \qquad q_j^-=\beta\frac{[h_jc_{uj}]^-}{\Sigma_u^-+\varepsilon}.
\end{align}
Bias and outside-option shares are absorbed. An empty branch is completed by a transition to
the cemetery, and its ordinary-successor fractions are zero ($0/0\coloneqq0$).
For unit-source hidden occupations $\Gamma^{\mathrm{RG}}_{j,\pm}(i)$, define
\begin{align}\label{ptnb:eq:affine-readout-occupation}
    \widehat\Gamma_\pm(i)=\sum_j\bigl(
    q_j^+\Gamma^{\mathrm{RG}}_{j,\pm}(i)
    +q_j^-\Gamma^{\mathrm{RG}}_{j,\mp}(i)\bigr).
\end{align}
Since $h_j=0$ implies $q_j^+=q_j^-=0$, only positive hidden activations receive nonzero relevance.
Subtracting the equations gives affine $\alpha\beta$ redistribution followed by
Theorem~\ref{ptnb:thm:ab-recovery}'s hidden-layer recursion. For any real initial relevance $r_u$,
\begin{align}\label{ptnb:eq:affine-readout-relevance}
    R_i=r_u\bigl(\widehat\Gamma_+(i)-\widehat\Gamma_-(i)\bigr).
\end{align}
The trajectory law and both occupations remain non-negative; the signed seed scales their
\emph{difference}. Logarithms apply only to non-negative branch masses, never the signed logit.
Classical LRP seeds $r_u=f_u(x)$; the released Routing-Game routine uses $r_u=1$, so the former map
is $f_u(x)$ times the latter, reversing sign for negative logits. At zero logits, output-seeded
relevance vanishes while unit-seeded relevance can remain nonzero.
